\section{The full-matrix distance and classical complement counts}
\label{sec:full-distance}

\subsection{Uniform expansion of the exact testing distance}

The nonsingular probability of a uniform alternating \(b\times b\)
matrix is \(\theta_{b/2}\). Apply
\eqref{eq:defect-expansion} at zero to obtain
\[
 \Pbb(D(U)\text{ nonsingular})
    =\theta_{b/2}+u\,\frac{q\theta_\infty}{q^2-1}
                  +O_q(uv).
\]
The tails of the infinite products give, uniformly for \(a\ge b\),
\begin{align}
 \theta_{b/2}
   &=\theta_\infty\left(1+\frac{qv}{q^2-1}
                                  +O_q(v^2)\right),\label{eq:theta-tail}\\
 \alpha_{a,b}
   &=\eta_\infty\left(1+\frac{u+v}{q^2-1}
                                  +O_q((u+v)^2)\right).\label{eq:alpha-tail}
\end{align}
For example,
\(\sum_{r>b/2}q^{-(2r-1)}=qv/(q^2-1)\);
the quadratic error follows by bounding the sum of all products of
two or more tail deficits. For \(\alpha_{a,b}\), the tail at
\((a+b)/2\) is already of order \(uv\).

It follows that
\begin{equation}\label{eq:defect-nonsing-asymp}
 \Pbb(D(U)\text{ nonsingular})
  =\theta_\infty
       \left(1+\frac{q(u+v)}{q^2-1}
                    +O_q((u+v)^2)\right).
\end{equation}
Insert \eqref{eq:alpha-tail} and \eqref{eq:defect-nonsing-asymp} into the exact
identity \eqref{eq:full-exact} to obtain
\[
 \TV(A,U)
 =\theta_\infty(1-\eta_\infty)
  +\frac{\theta_\infty\{q-(q+1)\eta_\infty\}}{q^2-1}
          (u+v)+O_q((u+v)^2).
\]
This proves \cref{thm:full}.
To check the sign, \(\eta_\infty>(q^2-2)/(q^2-1)\), while
\[
 \frac{q^2-2}{q^2-1}-\frac q{q+1}
     =\frac{q-2}{q^2-1}\ge0.
\]
The strict product inequality handles \(q=2\), so
\(\eta_\infty>q/(q+1)\) and \(c_q<0\).

\subsection{An independent positive-sum formula}

De Boeck and Van de Voorde's complement count \cite[Theorem~4.9]{DBV}
specializes to the following convenient formula. With \(h=b/2\), set
\begin{equation}\label{eq:positive-sum}
 p_{a,b}
 =\sum_{m=0}^{h}\theta_{h-m}
          \frac{\eta_h}{\eta_m\eta_{h-m}}\,
                    q^{-(a+1)m}.
\end{equation}
Then
\begin{equation}\label{eq:positive-sum-meaning}
 p_{a,b}=\Pbb(D(U)\text{ nonsingular}).
\end{equation}
Here is the identification of their count with the present model.
A nondegenerate \(b\)-subspace of \(V\) is disjoint from the target
\(a\)-summand exactly when the hidden square block \(B\) in
\eqref{eq:gram-condition} is invertible.
Since \(D(A)=B^{\trans}J_bB\), this is equivalent to
\(D(A)\) being nonsingular.
On this event the likelihood is the constant \(\alpha_{a,b}\).
If \(\rho\) denotes their disjointness probability, then
\[
 \rho=\Pbb(B\text{ invertible})
       =\alpha_{a,b}\Pbb(D(U)\text{ nonsingular}).
\]
After their finite products and powers of \(q\) are normalized,
Theorem~4.9 gives exactly
\(\rho/\alpha_{a,b}=p_{a,b}\).
Thus the finite scalar formula \eqref{eq:positive-sum} is a
reformulation of a published count, not a new enumeration theorem.

Formula \eqref{eq:positive-sum} independently confirms
\eqref{eq:defect-nonsing-asymp}. Its \(m=0\) term is
\(\theta_h\). Its \(m=1\) term is
\[
 \theta_{h-1}
    \frac{1-q^{-2h}}{1-q^{-2}}\,q^{-a-1}
       =\frac{q\theta_\infty}{q^2-1}\,u+O_q(uv).
\]
All terms with \(m\ge2\) sum to \(O_q(u^2)\), since
\(\eta_h/(\eta_m\eta_{h-m})\le\eta_\infty^{-2}\)
and their powers form a dominated geometric series.
This supplies a second derivation of the full-matrix expansion.
It also gives a stable positive-term formula for numerical evaluation,
independent of alternating-moment inversion.

\begin{remark}[The requirement that both dimensions grow]\label{rem:fixed-b}
If \(b\) is held fixed and \(a\to\infty\), the exact formulas instead give
\[
 \TV(A,U)\longrightarrow
        \theta_{b/2}(1-\eta_{b/2}).
\]
For instance, at \(q=2,b=2\) this limit is \(1/8\), not \(\Delta_2\).
Thus the universal constant in \cref{thm:full} requires both dimensions
to grow. The uniform signed estimates themselves hold at every
positive even size.
\end{remark}

